% Slide build of Figure 22.1 -- the four NAT names, one at a time.
%
% This is the figure in Part IV that most repays being built. Students do not
% find NAT hard; they find "inside local / inside global / outside local /
% outside global" hard, and they find it hard because all four arrive at once
% and look like four addresses. They are two addresses under four names.
%
% So each step adds exactly one label and says which half of the name it just
% fixed: inside/outside says whose host it is, local/global says who is
% looking. By step 4 the room has been told the rule four times without the
% lecturer repeating it once.
\begin{tikzpicture}[font=\scriptsize]
  \useasboundingbox (-2.0,-4.6) rectangle (12.6,2.0);

  \tikzset{
    dv/.style={rounded corners=3pt, draw=#1, fill=#1!8, text=#1!45!black,
               line width=0.9pt, minimum width=1.45cm, minimum height=1.05cm,
               align=center, font=\scriptsize\bfseries},
    tag/.style={font=\tiny, align=center},
    stepcap/.style={font=\bfseries, anchor=north west, text width=14.2cm,
                    align=left}}

  \node[dv=neutral] (pc)  at (0,0)    {\faIcon{desktop}\\PC1};
  \node[dv=dom3]    (r)   at (5.3,0)  {\faIcon{route}\\R1};
  \node[dv=dom1]    (web) at (10.6,0) {\faIcon{server}\\Server};

  \draw[line width=0.9pt, neutral!75] (pc) -- (r);
  \draw[line width=0.9pt, neutral!75] (r) -- (web);

  \draw[dom4!50, dash pattern=on 3pt off 2pt, line width=0.8pt]
        (5.3,1.5) -- (5.3,-2.3);
  \node[font=\tiny\bfseries, text=dom4!45!black] at (3.2,1.35) {INSIDE};
  \node[font=\tiny\bfseries, text=dom4!45!black] at (7.4,1.35) {OUTSIDE};
  \node[font=\tiny, text=neutral!85!black] at (2.4,0.55) {\cmd{ip nat inside}};
  \node[font=\tiny, text=neutral!85!black] at (8.2,0.55) {\cmd{ip nat outside}};

  \stepvis{1}{\node[tag, text=dom4!45!black] at (0,-1.15)
      {\textbf{inside local}\\192.168.10.5};}
  \stepvis{2}{\node[tag, text=dom4!45!black] at (3.4,-1.9)
      {\textbf{inside global}\\203.0.113.2};}
  \stepvis{3}{\node[tag, text=dom1!45!black] at (10.6,-1.15)
      {\textbf{outside global}\\198.51.100.7};}
  \stepvis{4}{\node[tag, text=dom1!45!black] at (7.2,-1.9)
      {\textbf{outside local}\\198.51.100.7};}

  % Step 4 is the moment to say out loud what the picture has been showing:
  % the server's two names are the same address, because nothing translates
  % it. That is why the fourth term feels pointless -- and it usually is.
  \stepvis{4}{\draw[{Circle[length=1.4mm]}-{Circle[length=1.4mm]},
                    dom1!55, line width=0.7pt,
                    dash pattern=on 2pt off 2pt] (8.45,-1.9) -- (9.55,-1.25)
      node[font=\tiny\itshape, text=dom1!45!black, midway, below=1pt,
           sloped] {identical};}

  \steponly{1}{\node[stepcap, text=dom4!45!black] at (-2.0,-3.2)
      {\emph{Inside} = it is our host. \emph{Local} = the address as the
       inside sees it. So: the real address on PC1.};}
  \steponly{2}{\node[stepcap, text=dom4!45!black] at (-2.0,-3.2)
      {Same host, other viewpoint. \emph{Global} = the address the outside
       world sees --- what R1 translated it to.};}
  \steponly{3}{\node[stepcap, text=dom1!45!black] at (-2.0,-3.2)
      {Now the other host. \emph{Outside} = theirs; \emph{global} = as
       their side sees it. The server's real address.};}
  \steponly{4}{\node[stepcap, text=dom1!45!black] at (-2.0,-3.2)
      {And the fourth is the same address again --- nothing translates the
       server, so outside local and outside global match. Two hosts, four
       names.};}
\end{tikzpicture}
